\subsection{Gripper timing does not clearly improve prediction of lost success}
\label{sec:supp-landmark-final}
Does the timing of the first opening-to-closing gripper command add information about lost success beyond the size of the command changes? At a cutoff of 60 controls, fixed in advance, adding timing lowers the held-out Brier loss slightly in both policies, with intervals that include zero (Table~\ref{tab:landmark-final}). This is an additional, exploratory study; the main paper does not rely on it. The feature is a command event, not physical closure or object contact.

\begin{table*}[t]
\caption{\textbf{Adding gripper timing to the command-change sizes does not clearly lower the loss at a fixed cutoff of 60 controls.} Cutoff pairs (of 404 per policy) reach the cutoff with valid command records and outcomes; the subset whose clean run succeeds (N/tasks) is analyzed. Improvement is baseline minus augmented out-of-fold Brier loss, positive when timing lowers the loss, with a 97.5\% paired whole-task interval. The denominator 404 is the fixed plan per policy; 375/360 pairs are analyzed.}

\centering\small
\begin{tabular}{lrrrrr}\toprule
Policy & Cutoff pairs & Clean-success N/tasks & Baseline Brier & Augmented Brier & Improvement [97.5\% interval]\\\midrule
$\pi_{0.5}$ & 396/404 & 375/27 & 0.09664 & 0.09555 & $0.00109[-0.00224,0.00448]$\\
$\pi_0$ & 397/404 & 360/26 & 0.12323 & 0.12172 & $0.00151[-0.00471,0.00724]$\\\bottomrule
\end{tabular}

\label{tab:landmark-final}
\end{table*}

\paragraph{Event reconstruction.}
Across all 50 runs, the Panda gripper command is positive for closing and negative for opening; zero leaves the accumulated command unchanged. Events use only complete, finite command records with valid identifiers up to the cutoff; the 39 unknown run outcomes, early terminations and technical failures are kept.

\paragraph{Baseline and augmented predictors.}
Of the 404 primary pairs per policy, both runs must reach the cutoff with valid command records and outcomes, and the clean run must succeed. Both predictors use only information available at the cutoff. The baseline uses $\log(1+\mathrm{RMS})$ of the translation, rotation and continuous-gripper command differences over the same first 60 controls, plus family indicators. The augmented predictor adds whether each run has an opening-to-closing transition by the cutoff, the difference of their times divided by 60 when both do, and each run's initial command direction; occurrence flags separate an absent event from a zero time difference.

\paragraph{Fitting on held-out tasks and intervals.}
The ridge model, fixed in advance, uses penalty $\alpha=1$, an unpenalized intercept, centering and scaling on the training data only, and leave-one-task-out prediction clipped to $[0,1]$; tasks are weighted equally, and families, settings and noise draws equally within their parents. Both predictors use identical rows in each fold; no training fold has a single class (which would fall back to the training mean). Both initial-closing indicator columns are zero in all eligible rows and add no variation. Paired whole-task resampling (2,000 draws) gives 97.5\% intervals per policy; the predictors are not refitted inside the bootstrap.

\paragraph{Outcomes of the analyzed pairs.}
On the primary support the outcome counts are $328/47/17/4$ SS/SF/FS/FF with eight unknown pairs for $\pi_{0.5}$ and $309/51/10/27$ with seven unknown pairs for $\pi_0$; every validated pair is kept whether or not a transition occurs. Pairs whose clean run succeeds number $93/98/82/102$ (Background/Camera/Light/Sensor Noise) for $\pi_{0.5}$ and $88/89/81/102$ for $\pi_0$. The estimates in Table~\ref{tab:landmark-final} apply to this conditional support.
